% Figure 1: one CONFLICT ladder (texts and scores from examples/teaser_<fid>.tex).
\begin{figure*}[t]
\centering
\begin{tikzpicture}[
  rung/.style={draw=lbline,rounded corners=1.2mm,fill=white,text width=4.45cm,
               inner sep=2.2mm,align=left,font=\scriptsize},
  head/.style={font=\footnotesize\bfseries,anchor=south west,inner sep=0pt},
  cl/.style={circle,draw=black!60,fill=white,inner sep=0pt,minimum size=4.6mm,
             font=\tiny},
  cnfedge/.style={draw=lbconflict,line width=0.9pt,dash pattern=on 2.2pt off 1.4pt},
  supedge/.style={draw=lbsupport,line width=0.9pt,-{Stealth[length=1.4mm]}},
  noedge/.style={draw=none},
]
\foreach \r/\C/\x/\title in {0/A/0/{two conflicts drawn},
                             1/B/5.3/{one conflict drawn},
                             2/C/10.6/{no conflict drawn}} {
  \node[rung,anchor=north west] (box\r) at (\x,0) {\parbox[t][3.7cm][t]{\linewidth}{%
    \slotfill{S1}{\csname teaser\C a\endcsname}\par\smallskip
    \slotfill{S2}{\csname teaser\C b\endcsname}}};
  \node[head] at ([yshift=1.2mm]box\r.north west) {$\rho_\r$: \title};
  % mini relation graph of the two slots
  \begin{scope}[shift={(\x,-4.6)}]
    \node[cl] (p\r) at (0.55,0) {$a_1$};
    \node[cl] (q\r) at (1.45,0) {$a_2$};
    \node[cl] (u\r) at (2.35,0) {$a_3$};
    \node[cl] (v\r) at (3.25,0) {$a_4$};
    \draw[\csname teaserEdge\C a\endcsname] (p\r) -- (q\r);
    \draw[\csname teaserEdge\C b\endcsname] (u\r) -- (v\r);
    % LCS bar
    \fill[lbline!40] (3.75,-0.11) rectangle (4.75,0.11);
    \fill[lbsupport] (3.75,-0.11) rectangle ({3.75+1.0*\csname teaserLcs\C\endcsname},0.11);
    \node[font=\tiny,anchor=east] at (4.75,0.3) {gold $\LCS=\csname teaserLcs\C\endcsname$};
  \end{scope}
}
\draw[-{Stealth[length=2mm]},lbline,line width=0.8pt]
  (0,-5.25) -- (15.55,-5.25)
  node[midway,below,font=\scriptsize,text=black!70] {declared: each rung more coherent than
  the one before (fewer violated conflicts at the committed world)};
\end{tikzpicture}
\caption{\textbf{A coherence ladder.} Family \teaserFamily{} answers ``\teaserQuestion''
in \teaserWords{} words. Its three rungs are the same answer with different text spliced into two
\emph{slots}; everything outside the highlighted sentences, every claim and every claim's truth
value is identical. In $\rho_0$ both slots set two true claims against each other
(``\ldots{} cannot be squared with \ldots''), in $\rho_1$ one does, in $\rho_2$ neither: under the
glyphs, a dashed edge is a drawn conflict between true claims. The gold relation graphs order the
rungs, and the gold $\LCS$ of the companion paper \citep{anoncoherence} rises accordingly. A
metric evaluated on LoCoBench must rank the rungs in the declared order.}
\label{fig:teaser}
\end{figure*}
